\section{Introduction}
\label{sec:introduction}

A bank that validates a probability-of-default model hands the validator one number to defend: the discrimination of the model on data it has not seen, usually the area under the receiver operating characteristic curve (AUC). Supervisory guidance on model risk management lists out-of-sample and out-of-time testing among the activities that test a model and asks for challenge by objective experts with enough independence to maintain objectivity \citep{sr262}; the instructions for reporting the validation of internal credit models ask for the discrimination of the current sample against the value observed at development \citep{ecb2019}. Both texts presuppose that the number on the unseen data is unseen in fact. The presupposition fails in a recurrent and silent way. A binning, a weight of evidence, an imputation median, a feature screen or a choice among configurations is fitted with rows that the estimate later treats as new, or a calibration map is fitted on the very rows that it is scored on, and the estimate drifts upward without any signal that a reviewer can read from the number itself \citep{kaufman2012,kapoor2023}. The literature on this failure is mature for microarray classification and for machine-learning science \citep{ambroise2002,varma2006,cawley2010,kapoor2023}; it is thin for the scorecards and challenger models that a credit risk function validates every year, where the preprocessing is target-aware by construction (weight of evidence, information value, supervised binning) and the data are often a single snapshot with no date on which to split.

The practical consequence is that the validation of a development team's model by an independent validator is only as good as the split. A validator who receives a model, a holdout and a report of the holdout AUC cannot tell from the AUC whether the holdout was read during fitting. What the validator can do is run the fitting path again with the outside information changed, and ask whether anything that was fitted moves. This is the idea of the protocol that this paper measures: the split is declared per case before any fit, three invariance checks (the gates) are executed on the fitting path and must pass, and the artifact that carries the results also carries the split, its sizes, its seed and the claim that the case makes or declines to make about the future.

The protocol is an enforcement of rules that the field already states \citep{hastie2009,siddiqi2017,kaufman2012}: transforms are fitted on the training slice only; a calibration map is fitted on a slice disjoint from the test slice; every fitted parameter is unchanged when a value outside the training information set is perturbed. It is not a new estimator. What this paper adds is the measurement that the rules lack. It states what each violation costs (the bias of the holdout AUC, the Brier score and the calibration error, and the coverage of the interval that the engine reports), what the three gates detect and what they miss, what a leak that exists only in development looks like to the tests a validation function already runs, and how far a random split inside the training window is from the out-of-time cohort that a dated case allows. Every number is the output of a script in this repository, run on the engine \texttt{riskvalidation} 1.1.0 and on the pipeline of the product that applies the protocol to its cases.

\paragraph{Contributions.} The paper makes the following contributions.
\begin{enumerate}
\item A protocol with a precise formal statement (Section~\ref{sec:proposed-method}): the row sets that each fitting step may read, the split declared per case, the rule that fixes the training window of a dated case, and three invariance gates with a non-vacuity control, written as an algorithm whose cost is a small multiple of one fit.
\item Five propositions and one conjecture (Section~\ref{sec:theory}): the holdout estimate is unbiased for the fitted model when no step read the holdout; the exact probability that the outside-information gate detects a leak through a quantile edge; a bound on the optimism of a choice made on the holdout; closed forms for the AUC of a score that contains a feature that exists only in development; the second-order form of the population stability index of that feature; and the detection of a discrete choice by repeated perturbation. The scaling of the bias of cross-validation with preprocessing outside the folds is stated as a conjecture with its evidence.
\item A measurement of the bias that each of nine leaks adds to the holdout estimate on planted data with a known truth, on 2,500-row samples of 30,000 card holders and on the product's own scorecard, with the coverage of the 95\% interval that the engine reports (Section~\ref{sec:results}).
\item A measurement of what the gates detect on the same leaks, with a subset ablation that shows which gate is necessary for which leak (Section~\ref{sec:ablation}).
\item A measurement, with the engine's own detectors, of what a leak that exists only in development looks like: the AUC-against-development test sees it, the stability index of the leaked feature and the single-feature screen do not.
\item A comparison, on the SBA 7(a) loans of five cohorts (training windows of \nv{e4.2006-07.window} to \nv{e4.2018-20.window} loans), of a random split inside the training window with the declared out-of-time cohort.
\end{enumerate}

\paragraph{Main results.} On planted data of \nv{syn.n} rows with 40 inputs, the holdout AUC of a WoE scorecard under the protocol has a bias of \nv{syn.w.none.bias}$\,\pm\,$\nv{syn.w.none.ci} (95\% interval over \nv{syn.reps} repetitions) against a reference sample of \nv{syn.oracle} rows that no step read, and the 95\% interval of the engine covers the reference value in \nv{syn.w.none.cov}\% of the repetitions. When the weight of evidence is computed with the holdout labels the bias is \nv{syn.w.woe.bias}$\,\pm\,$\nv{syn.w.woe.ci} and the coverage falls to \nv{syn.w.woe.cov}\%; when every step reads the holdout the bias is \nv{syn.w.fit.bias}$\,\pm\,$\nv{syn.w.fit.ci} and the coverage is \nv{syn.w.fit.cov}\%. At least one of the three gates fails for each of the \nv{h5.total} leaks planted in the fitting path in every one of \nv{e2.reps} draws, and none fails on the protocol itself. In the experimental design, \nv{h.notheld.word} of nine predictions did not hold: on the card holders the screen with holdout labels does not raise the AUC by more than the noise (H2); the out-of-time gap of the gradient-boosting rung is below 2 points in two of the five cohorts (H8); and the cross-validated AUC of the logit is not inflated on the card holders before the folds and is slightly negative inside the folds on the planted data (H9) (Section~\ref{sec:results}). The predictions were written together with the outputs that decide them, and H6 and H8 were set with committed values of the cases in view (Section~\ref{sec:experimental-design}). The leak that exists only in development inflates the development AUC by \nv{e3.s1.0.infl} at a shift of one standard deviation; the AUC-against-development test rejects in \nv{e3.s1.0.auc_test}\% of the repetitions while the stability index of the leaked feature rejects in \nv{e3.s1.0.csi_test}\% and the single-feature screen flags it in \nv{e3.s1.0.single_feature_flag}\%.

\paragraph{Scope.} The protocol is enforcement, not discovery. It guarantees nothing about a leak that the perturbations do not reach (a feature that is constructed from information dated after the observation date and that is already in the data is outside what the gates perturb), and Section~\ref{sec:limitations} states the threats to validity of the measurements that follow. The paper separates what the gates prove (the fitted parameters do not depend on the perturbed rows) from what they leave open (whether the training information set itself is free of the future).
